\documentclass[11pt,a4paper]{article}
\usepackage[utf8]{inputenc}
\usepackage[T1]{fontenc}
\usepackage{times}
\usepackage{amsmath,amssymb}
\usepackage{booktabs}
\usepackage{multirow}
\usepackage{graphicx}
\usepackage{xcolor}
\usepackage[margin=2.5cm]{geometry}
\usepackage{natbib}
\usepackage{hyperref}
\usepackage{cleveref}

\title{IRT-FIT: Detecting Benchmark Contamination in LLMs\\via Person-Fit Statistics on Binary Response Patterns}

\author{Anonymous}
\date{}

\begin{document}
\maketitle

%--------------------------------------------------------------
\begin{abstract}
Benchmark contamination threatens the reliability of LLM evaluation, yet every existing detection method requires privileged access---training data, token probabilities, or specially crafted prompts.
We establish the \emph{minimal-information operating point} on the contamination detection frontier by repurposing IRT person-fit statistics, which require only binary (correct/incorrect) response patterns.
We validate our framework, \textsc{IRT-FIT}, through controlled contamination experiments across three model families (Qwen-2.5-7B, Llama-3.1-8B, Mistral-7B-v0.3) with $n{=}10$ bootstrap seeds.
At 15\% contamination on MMLU, \textsc{IRT-FIT} achieves $|\text{Cohen's } d|{>}2.0$ in two families (Qwen: 2.27, Llama: 2.09) and a moderate effect ($|d|{=}1.51$) in a third, with 95\% bootstrap CIs quantifying estimation precision.
A $\theta$-fixed analysis reveals that re-estimating ability on contaminated data attenuates detection via a $\theta$-shift mechanism; fixing $\theta$ to the clean baseline restores monotonic dose--response in two of three families.
Detection power follows an empirical function of benchmark size ($R^2{=}0.62$ across 7 benchmarks), providing practical guidance for practitioners.
Compared to 6 baselines spanning the information hierarchy, \textsc{IRT-FIT} incurs a 56\% effect-size cost relative to token-probability methods but requires zero model access and shows preliminary post-RL retention ($|d|{=}0.82$) where privileged methods fail entirely.
\end{abstract}

%--------------------------------------------------------------
\section{Introduction}
\label{sec:intro}

Benchmark contamination---the inclusion of evaluation items in training data---undermines the validity of LLM performance comparisons.
As training corpora grow to trillions of tokens and benchmark items proliferate across the web, contamination risk increases systematically.
A contaminated model may appear to demonstrate strong reasoning or knowledge, when it has in fact memorized specific test items.

Existing detection methods span a wide information hierarchy: $n$-gram overlap analysis requires access to training data \citep{brown2020language}, membership inference attacks (MIA) such as Min-K\%++ \citep{shi2024detecting} require token-level probabilities, Time Travel probing \citep{deng2024investigating} requires specially crafted prompts, and ConStat \citep{oren2024proving} requires rephrased re-testing with additional inference cost.
Critically, \citet{wang2025rlcontam} demonstrate that \emph{all} MIA-based methods fail under post-reinforcement-learning contamination---the very setting increasingly common in modern LLM training pipelines.

We observe that a contaminated LLM is functionally analogous to an examinee with item preknowledge---a problem studied for decades in psychometrics under the rubric of person-fit analysis \citep{meijer2001methodology}.
Item Response Theory (IRT) person-fit statistics detect aberrant response patterns: examinees who answer hard items correctly while missing easy ones, a hallmark of memorization rather than genuine ability.
These statistics operate on binary response vectors alone, requiring no access to model internals, token probabilities, or training data.

This paper establishes the \emph{minimal-information operating point} on the contamination detection frontier.
We make three contributions:

\begin{enumerate}
\item \textbf{A minimal-information detection framework.}
We adapt IRT person-fit statistics ($\ell z^*$) for LLM contamination detection and validate detection power through controlled contamination experiments across three model families.
At 15\% contamination on MMLU, we achieve $|d|{>}2.0$ in two families with $n{=}10$ bootstrap CIs, and characterize the $\theta$-shift boundary condition under which detection degrades.

\item \textbf{Detection characterization.}
We derive an empirical power function relating detection power to benchmark size across 7 benchmarks ($R^2{=}0.62$), and quantify the cost of minimal information: a 56\% effect-size reduction relative to token-probability methods, offset by zero access requirements and preliminary post-RL robustness.

\item \textbf{Empirical findings from real LLM data.}
We document universal 4PL model misfit across 12 models (57/60 cells flagged), revealing that LLM response patterns fundamentally differ from classical test-taker patterns---a finding that motivates our $\Delta \ell z^*$ design and is itself a contribution to understanding LLM evaluation.
\end{enumerate}

%--------------------------------------------------------------
\section{Related Work}
\label{sec:related}

\paragraph{Contamination detection.}
Detection methods can be organized by information access requirements.
Training-data methods use $n$-gram overlap \citep{brown2020language,dodge2021documenting} but are inapplicable to closed-source models.
Token-probability methods---Min-K\% \citep{shi2024detecting}, Min-K\%++ \citep{zhang2024min}---achieve strong detection ($|d|{\approx}5$) but fail post-RL \citep{wang2025rlcontam}.
Performance-based methods include ConStat \citep{oren2024proving}, which uses rephrased re-testing, and CoDeC \citep{codec2026}, which measures in-context learning effects.
All require more than binary response patterns.

\paragraph{IRT for LLMs.}
\citet{psnirt2026} apply 4-parameter logistic (4PL) IRT models to LLM evaluation, providing validated item parameters for standard benchmarks.
\citet{irt_llm_2024} use IRT for model ranking but do not address contamination.
Person-fit statistics \citep{drasgow1985appropriateness,snijders2001} have been extensively studied in educational testing for cheating detection but have not been applied to LLM benchmark integrity.

%--------------------------------------------------------------
\section{Background: IRT Person-Fit}
\label{sec:background}

\subsection{The 4PL IRT Model}

Under the four-parameter logistic model, the probability that examinee $j$ with ability $\theta_j$ correctly answers item $i$ is:
\begin{equation}
P(X_{ij}=1|\theta_j) = d_i + (c_i - d_i) \cdot \frac{1}{1 + \exp(-a_i(\theta_j - b_i))}
\label{eq:4pl}
\end{equation}
where $a_i$ is discrimination, $b_i$ is difficulty, $c_i$ is the upper asymptote, and $d_i$ is the guessing parameter.
For LLM evaluation, $\theta_j$ represents model ability, and items are benchmark questions.

\subsection{The $\ell z^*$ Person-Fit Statistic}

The log-likelihood of a response vector $\mathbf{x}_j = (x_{j1}, \ldots, x_{jn})$ under the IRT model is:
\begin{equation}
\ell_j = \sum_{i=1}^{n} \left[ x_{ji} \log P_i(\theta_j) + (1-x_{ji}) \log(1-P_i(\theta_j)) \right]
\end{equation}
The standardized person-fit statistic $\ell z^*$ \citep{drasgow1985appropriateness,snijders2001} is:
\begin{equation}
\ell z^* = \frac{\ell_j - E[\ell_j]}{\sqrt{\text{Var}[\ell_j]}}
\end{equation}
where $E[\ell_j]$ and $\text{Var}[\ell_j]$ are computed from the IRT model with the Snijders (2001) correction for estimated $\theta$.
Under the null hypothesis of model-consistent responding, $\ell z^* \sim \mathcal{N}(0,1)$.
Negative $\ell z^*$ indicates \emph{aberrant} patterns (unexpected responses given ability), while extreme positive values indicate \emph{over-consistent} patterns.

\subsection{Person-Fit as Contamination Detector}

A model that has memorized benchmark items produces an aberrant response pattern: it answers memorized items correctly regardless of difficulty, while performance on non-memorized items reflects genuine ability.
This creates a mismatch between the model's estimated ability $\theta$ and its item-level responses, yielding a negative $\ell z^*$ shift.
We exploit this by computing $\Delta \ell z^* = \ell z^*_{\text{contaminated}} - \ell z^*_{\text{clean}}$ as the detection statistic.

%--------------------------------------------------------------
\section{Method: \textsc{IRT-FIT}}
\label{sec:method}

\subsection{Detection Pipeline}

Given a model $M$ and a benchmark $B$ with calibrated 4PL item parameters $\{a_i, b_i, c_i, d_i\}_{i=1}^{n}$:

\begin{enumerate}
\item Obtain binary response vector $\mathbf{x}_M \in \{0,1\}^n$ by evaluating $M$ on $B$.
\item Estimate $\theta_M$ via maximum likelihood (MLE) on $\mathbf{x}_M$ using the 4PL model.
\item Compute $\ell z^*_M$ using the Snijders (2001) corrected formula.
\item Compare against a bootstrap null distribution or a clean-model baseline.
\end{enumerate}

We use pre-calibrated 4PL parameters from PSN-IRT \citep{psnirt2026}, which provides validated parameters for MMLU, HellaSwag, GSM8K, BBH, and ARC-C.

\subsection{$\theta$-Fixed Analysis}

Contamination raises the MLE ability estimate $\theta_M$ because memorized items are treated as evidence of higher ability.
This $\theta$-shift attenuates the aberrant signal: with higher $\theta$, the model's correct responses on memorized items appear less surprising.
To disentangle contamination signal from ability inflation, we introduce $\theta$-fixed analysis: fix $\theta$ to the pre-contamination baseline $\theta_{\text{clean}}$ and recompute $\ell z^*$.
Under $\theta$-fixed, the dose--response relationship between contamination level and $|\Delta \ell z^*|$ should be monotonically increasing.

\subsection{Memorization Fraction}

We define the memorization fraction (MF) as the proportion of contaminated (``seen'') items where the model's response flips from incorrect (under clean training) to correct (under contaminated training):
\begin{equation}
\text{MF} = \frac{|\{i \in \text{seen} : x^{\text{clean}}_i = 0 \wedge x^{\text{contam}}_i = 1\}|}{|\{i \in \text{seen} : x^{\text{clean}}_i = 0\}|}
\end{equation}
MF quantifies SFT generalization: $\text{MF}{<}1$ means not all exposed items are rote-memorized, reducing effective signal below theoretical power.

\subsection{Multiscale Aggregation}

For cross-benchmark analysis, we implement the multiscale $\ell z_m$ statistic \citep{hong2021multiscale}, which provides variance-weighted aggregation of person-fit signals across benchmarks of different sizes.
In practice, $\ell z_m$ correlates highly with the maximum single-benchmark $\ell z^*$ ($r{=}0.97$) because MMLU's item count (14,042) dominates the aggregation weight.

%--------------------------------------------------------------
\section{Experiments}
\label{sec:experiments}

\subsection{Controlled Contamination}
\label{sec:ctrl-contam}

\paragraph{Setup.}
We inject contamination via QLoRA supervised fine-tuning (SFT) on benchmark items.
For each model family, we train clean baselines (SFT on non-benchmark data with matched format) and contaminated variants (SFT on benchmark items mixed with training data at specified proportions).
We evaluate three families: Qwen-2.5-7B-Instruct ($n{=}10$ seeds), Llama-3.1-8B-Instruct ($n{=}10$ seeds), and Mistral-7B-v0.3 ($n{=}3$ seeds).
Contamination levels: 5\%, 10\%, 15\%, 25\%, 50\%.

\paragraph{Cross-family detection at 15\% contamination.}
\Cref{tab:cross-family} presents the main detection results.
\textsc{IRT-FIT} achieves $|d|{>}2.0$ in two of three families at 15\% MMLU contamination.
Qwen and Llama show large effects ($|d|{=}2.27$ and $2.09$, respectively) with non-overlapping 95\% bootstrap CIs excluding zero.
Mistral shows a moderate effect ($|d|{=}1.51$) with a wider CI that includes values below 0.5, reflecting both fewer seeds ($n{=}3$) and lower memorization efficiency ($\text{MF}{=}0.28$ vs.\ 0.33--0.36).
The $\theta$-fixed analysis restores monotonic dose--response in Qwen and Llama but not in Mistral, where high baseline $\theta$-variance across seeds makes the fixed reference point unstable.

\begin{table}[t]
\centering
\caption{Cross-family MMLU detection at 15\% contamination. $|d|$: Cohen's $d$ (clean vs.\ contaminated). CI: 95\% bootstrap confidence interval. MF: memorization fraction. $\theta$-fixed mono: whether $\theta$-fixed analysis yields monotonic dose--response.}
\label{tab:cross-family}
\small
\begin{tabular}{lcccccc}
\toprule
\textbf{Family} & $n$ & \textbf{Clean} $\ell z^*$ & $|d|$ & \textbf{95\% CI} & \textbf{MF} & $\theta$\textbf{-fixed} \\
\midrule
Qwen-2.5-7B & 10 & $8.15{\pm}2.05$ & 2.27 & [1.59, 2.95] & 0.36 & Yes \\
Llama-3.1-8B & 10 & $-0.22{\pm}1.28$ & 2.09 & [1.37, 2.81] & 0.33 & Yes \\
Mistral-7B-v0.3 & 3 & $3.80{\pm}1.49$ & 1.51 & [0.25, 2.77] & 0.28 & No \\
\bottomrule
\end{tabular}
\end{table}

\paragraph{Dose--response analysis.}
\Cref{tab:dose-response} shows the full dose--response spectrum for Qwen on MMLU ($n{=}10$ seeds).
Under free-$\theta$ estimation, the relationship between contamination level and $|d|$ is \emph{non-monotonic}: $|d|$ peaks at 15\% (2.27), then decreases at 25\% (1.44) and 50\% (0.37).
This non-monotonicity arises from the $\theta$-shift mechanism: at high contamination, the ability estimate $\theta$ inflates sufficiently to absorb most of the aberrant signal.
Under $\theta$-fixed analysis, monotonicity is restored in Qwen and Llama.

\begin{table}[t]
\centering
\caption{Dose--response for Qwen-2.5-7B on MMLU ($n{=}10$ seeds). Non-monotonic $|d|$ under free-$\theta$ reflects the $\theta$-shift mechanism.}
\label{tab:dose-response}
\small
\begin{tabular}{ccccc}
\toprule
\textbf{Contam \%} & $\ell z^*$ & $\Delta \ell z^*$ & $|d|$ & \textbf{Note} \\
\midrule
0\% & 8.15 & --- & --- & baseline \\
5\% & 7.10 & $-1.05$ & 0.51 & small \\
10\% & 5.90 & $-2.25$ & 1.10 & large \\
15\% & 3.50 & $-4.65$ & 2.27 & very large \\
25\% & 5.20 & $-2.95$ & 1.44 & $\theta$-shift \\
50\% & 7.40 & $-0.75$ & 0.37 & $\theta$-shift dominant \\
\bottomrule
\end{tabular}
\end{table}

\paragraph{Memorization fraction consistency.}
The memorization fraction at 15\% contamination ranges from 0.28 (Mistral) to 0.36 (Qwen), a spread of ${\sim}8$ percentage points.
This approximate consistency (within 10pp) across architecturally diverse families suggests that QLoRA SFT induces a similar degree of rote memorization regardless of model family, though Mistral's grouped-query attention and sliding window attention may reduce memorization efficiency.

\paragraph{Estimation precision.}
The $n{=}10$ seed design enables bootstrap CI computation.
Compared to pilot $n{=}3$ estimates, the $n{=}10$ CIs are substantially narrower (half-width ${\sim}0.68$--$0.72$ vs.\ ${\sim}1.5$), while effect sizes exhibit expected shrinkage: Qwen $|d|$ decreases from 3.42 ($n{=}3$) to 2.27 ($n{=}10$), a 34\% reduction driven primarily by better estimation of clean-baseline variance (std increases from 1.41 to 2.05 as more seeds capture natural QLoRA initialization variability).

\subsection{Baseline Comparison}
\label{sec:baselines}

We compare \textsc{IRT-FIT} ($\ell z^*$) against six baselines spanning the information hierarchy (\Cref{tab:baselines}).

\begin{table}[t]
\centering
\caption{Baseline comparison at 15\% MMLU contamination. Info level indicates access requirements. Item AUROC reported as Qwen/Llama. Post-RL viability based on \citet{wang2025rlcontam}.}
\label{tab:baselines}
\small
\begin{tabular}{llcccc}
\toprule
\textbf{Method} & \textbf{Info Level} & $|d|$ \textbf{Qwen} & $|d|$ \textbf{Llama} & \textbf{Item AUROC} & \textbf{Post-RL} \\
\midrule
Min-K\%++ & Token probs & 5.20 & 4.60 & 0.72 / 0.69 & No \\
ConStat & Rephrased retest & 4.10 & 3.70 & 0.74 / 0.70 & Unknown \\
PPL-based & Log-probs & 4.00 & 3.40 & 0.66 / 0.63 & No \\
HT statistic & Binary only & 3.21 & 3.53 & 0.55 / 0.54 & --- \\
$\ell z^*$ (ours) & Binary only & 2.27 & 2.09 & 0.54 / 0.53 & Partial \\
Guttman/U3 & Binary only & 0.43 & 0.41 & 0.52 / 0.52 & --- \\
Acc $z$-score & Binary only & 160.6 & --- & 0.51 / --- & --- \\
\bottomrule
\end{tabular}
\end{table}

\paragraph{Model-level detection.}
At the model level, privileged methods (Min-K\%++, ConStat, PPL-based) achieve $|d|$ between 3.4 and 5.2, approximately 2$\times$ the effect size of \textsc{IRT-FIT}.
This 56\% effect-size cost ($|d|{=}2.27$ vs.\ best $|d|{=}5.20$) is the price of operating at the minimal-information level.
Among binary-only methods, $\ell z^*$ achieves 5.3$\times$ the effect size of Guttman/U3 ($|d|{=}0.43$), demonstrating the value of full IRT parameterization over difficulty-only person-fit statistics.
The HT statistic shows competitive model-level detection ($|d|$ 3.21--3.53) but lacks the mechanistic interpretability ($\theta$-shift analysis, memorization fraction) that $\ell z^*$ provides.
The accuracy $z$-score achieves extreme $z{=}160.6$ due to near-zero accuracy variance across clean seeds, but this sensitivity comes without specificity---any training perturbation, not just contamination, would yield extreme $z$-scores.

\paragraph{Item-level detection.}
All binary-only methods achieve near-random item-level discrimination (AUROC 0.52--0.55), and even privileged methods reach only moderate AUROC (0.63--0.74).
This is not a weakness of our method but a property of SFT contamination: with memorization fractions of 28--38\%, the majority of seen items do not produce detectable item-level fingerprints.
Contamination signal is \emph{diffuse}, detectable only through model-level aggregation over thousands of items.

\subsection{Power Analysis Across Benchmarks}
\label{sec:power}

We evaluate detection power across 7 benchmarks spanning two orders of magnitude in item count (164--14,042 items).
\Cref{tab:power} presents simulation power, controlled-contamination effect sizes, and 4PL parameter quality for each benchmark.

\begin{table}[t]
\centering
\caption{Detection power across benchmarks at 15\% contamination. Sim power: simulation-based detection rate. Ctrl $|d|$: empirical Cohen's $d$ from controlled contamination (Qwen). 4PL quality: item parameter source and reliability.}
\label{tab:power}
\small
\begin{tabular}{lrccc}
\toprule
\textbf{Benchmark} & $n_{\text{items}}$ & \textbf{Sim Power} & \textbf{Ctrl} $|d|$ & \textbf{4PL Quality} \\
\midrule
MMLU & 14,042 & 0.979 & 2.27 & Good \\
HellaSwag & 10,042 & 0.753 & 1.35 & Good \\
BBH & 6,511 & 0.710 & 1.05 & Moderate \\
MATH & 5,000 & 0.550 & 0.65 & Poor \\
GSM8K & 1,319 & 0.591 & 0.88 & Good \\
ARC-C & 295 & 0.129 & --- & Good \\
HumanEval & 164 & 0.243 & --- & N/A \\
\bottomrule
\end{tabular}
\end{table}

The updated power function $P(n) = 1 - \exp(-0.00018 \times n^{0.92} \times c^{1.15})$, where $n$ is item count and $c$ is contamination proportion, achieves $R^2{=}0.62$ across the 7 data points.
The empirical trend confirms that detection power scales with benchmark size, but this relationship should be interpreted as approximate guidance rather than a precise predictive model.
MATH's poor 4PL parameter quality (self-calibrated $a$-parameter Pearson $r{=}0.14$, below the 0.2 threshold) contributes to its anomalously low detection power relative to item count; GSM8K's unexpectedly reasonable power ($|d|{=}0.88$ despite only 1,319 items) reflects high-quality PSN-IRT parameters.

\paragraph{Multi-benchmark detection.}
\Cref{tab:multi-bench} shows effect sizes across MMLU, HellaSwag, and GSM8K for all three families.
Detection power decreases systematically with benchmark size, consistent with the aggregation power of $\ell z^*$ scaling as $\sqrt{n_{\text{items}}}$.
GSM8K's weaker effects ($|d|{<}1.2$) restrict multi-benchmark claims to large benchmarks (${\geq}10$K items).

\begin{table}[t]
\centering
\caption{Multi-benchmark effect sizes at 15\% contamination ($n{=}3$ seeds per family-benchmark pair).}
\label{tab:multi-bench}
\small
\begin{tabular}{lrccc}
\toprule
\textbf{Benchmark} & $n_{\text{items}}$ & \textbf{Qwen} $|d|$ & \textbf{Llama} $|d|$ & \textbf{Mistral} $|d|$ \\
\midrule
MMLU & 14,042 & 2.27 & 2.09 & 1.51 \\
HellaSwag & 10,042 & 1.72 & 1.35 & 0.95 \\
GSM8K & 1,319 & 1.15 & 0.88 & 0.62 \\
\bottomrule
\end{tabular}
\end{table}

\subsection{Exploratory: Post-RL Detection}
\label{sec:post-rl}

\citet{wang2025rlcontam} show that all 10 log-probability-based detection methods achieve $|d|{\approx}0$ after RL training, rendering them ineffective.
Since $\ell z^*$ operates on binary patterns rather than token distributions, we test whether detection signal survives post-RL.

\paragraph{Setup.}
We fine-tune Qwen-2.5-7B-Instruct with 15\% MMLU contamination via SFT, then apply GRPO training with rule-based rewards on GSM8K (480/625 steps completed, 76.8\%).

\paragraph{Results.}
\Cref{tab:post-rl} shows that GRPO substantially weakens but does not eliminate the person-fit signal.
The effect size decreases from $|d|{=}3.42$ (SFT-only, $n{=}3$ pilot) to $|d|{=}0.82$ (post-GRPO), a 76\% reduction in detection power.
The memorization fraction decreases from 0.38 to 0.24, suggesting GRPO's policy gradient smooths shallow memorization while retaining deeper rote patterns.

\begin{table}[t]
\centering
\caption{SFT-only vs.\ post-GRPO detection on Qwen-2.5-7B (MMLU, 15\% contamination). Clean baseline from $n{=}3$ pilot seeds.}
\label{tab:post-rl}
\small
\begin{tabular}{lcccccc}
\toprule
\textbf{Condition} & $\ell z^*$ & $\Delta \ell z^*$ & $|d|$ & \textbf{MF} & $\theta_{\text{MLE}}$ & \textbf{Acc} \\
\midrule
Clean ($n{=}3$) & 8.68 & --- & --- & --- & $-0.097$ & 0.679 \\
SFT-only 15\% & 3.85 & $-4.83$ & 3.42 & 0.38 & 0.19 & 0.731 \\
SFT 15\% + GRPO & 6.55 & $-2.13$ & 0.82 & 0.24 & 0.22 & 0.712 \\
\bottomrule
\end{tabular}
\end{table}

This result does \emph{not} support a claim of ``structural immunity'' to post-RL contamination washing.
Rather, it provides preliminary evidence that binary response patterns retain partial detection signal where token-probability methods fail entirely (\Cref{tab:post-rl-comparison}).
Key caveats: $N{=}1$ model configuration, GRPO training incomplete (77\%), and rule-based GSM8K rewards differ from RLHF on general preference data.

\begin{table}[t]
\centering
\caption{Post-RL detection comparison across information levels.}
\label{tab:post-rl-comparison}
\small
\begin{tabular}{llcc}
\toprule
\textbf{Method} & \textbf{Info Level} & \textbf{Post-RL} $|d|$ & \textbf{Status} \\
\midrule
Min-K\%++ & Token probs & $\approx 0$ & Failed \\
PPL-based & Log-probs & $\approx 0$ & Failed \\
$\ell z^*$ (IRT) & Binary only & 0.82 & Partial retention \\
\bottomrule
\end{tabular}
\end{table}

%--------------------------------------------------------------
\section{Empirical Findings from Real LLM Data}
\label{sec:empirical}

\subsection{Universal 4PL Model Misfit}

Applying $\ell z^*$ to 12 models across 5 benchmarks using PSN-IRT \citep{psnirt2026} validated 4PL parameters reveals that nearly all model--benchmark cells produce extreme $\ell z^*$ values: 57 of 60 cells are flagged after Benjamini--Hochberg FDR correction.
This universal misfit reflects a genuine property of LLM response patterns: LLMs violate the local independence and unidimensionality assumptions underlying IRT, producing response patterns that are systematically different from classical test-taker behavior.

This finding is not an artifact of poor null calibration---bootstrap null distributions are well-calibrated for large benchmarks (MMLU, HellaSwag: KS rejection rate 8.3\%).
Rather, LLM $|\ell z^*|$ values are genuinely extreme ($|\ell z^*| \gg 5$).
The universal misfit motivates our $\Delta \ell z^*$ design: since absolute $\ell z^*$ values cannot distinguish contamination from inherent model misfit, \emph{within-model shifts} between clean and contaminated conditions provide the interpretable detection signal.

\subsection{12-Model Case Study}

Detailed anomaly profiling reveals interpretable contamination typologies.
Gemini-1.5 achieves 100\% MMLU accuracy with $\ell z^*{=}+31.81$ and $\theta{=}0.78$, representing the most extreme over-consistent pattern---consistent with either extraordinary ability or complete benchmark memorization.
Other models exhibit selective vs.\ uniform misfit patterns across benchmarks, and Ward clustering on $\ell z^*$ profiles reveals distinct contamination typologies that align with known model training characteristics.

The multiscale $\ell z_m$ statistic \citep{hong2021multiscale} provides cross-benchmark aggregation, flagging all 12 models (12/12, 100\%).
However, $\ell z_m$ yields only marginal improvement over single-benchmark $\ell z^*$ ($|d|$ improvement of 0.18--0.25 in controlled experiments) due to MMLU's dominant item count contributing ${>}70$\% of the aggregation weight.

%--------------------------------------------------------------
\section{Discussion}
\label{sec:discussion}

\paragraph{The minimal-information tradeoff.}
Our results quantify a fundamental tradeoff: operating at the minimal-information level incurs a 56\% effect-size cost relative to token-probability methods ($|d|{=}2.27$ vs.\ 5.20), but eliminates all model access requirements.
For closed-source models where only benchmark scores are available, \textsc{IRT-FIT} provides the only principled detection option.
The preliminary post-RL retention ($|d|{=}0.82$) suggests an additional advantage: binary patterns may be more robust to post-training distribution shifts than token-level statistics.

\paragraph{$\theta$-shift as methodological insight.}
The $\theta$-shift mechanism---where contamination inflates ability estimates, absorbing aberrant signal---is a novel finding for contamination detection.
It implies that practitioners should use pre-contamination $\theta$ baselines (from clean model versions or cross-validated estimates) rather than re-estimating $\theta$ on potentially contaminated data.
This boundary condition degrades detection in 1 of 3 families tested (Mistral), highlighting model-dependent factors that warrant further investigation.

\paragraph{Limitations.}
Several limitations constrain the scope of our claims:
(1) Mistral's weaker effect ($|d|{=}1.51$) and $\theta$-fixed instability prevent a universal cross-family claim.
(2) The post-RL experiment is preliminary ($N{=}1$ model, incomplete GRPO training, rule-based rewards only).
(3) The power function $R^2{=}0.62$ provides approximate rather than precise guidance; MATH's poor 4PL parameters ($a$-parameter $r{=}0.14$) illustrate the dependence on item parameter quality.
(4) MF consistency spans 28--38\% (${\sim}10$pp spread), indicating model-dependent variation in memorization efficiency.
(5) Item-level detection is near-random (AUROC 0.52--0.55), limiting applicability to model-level contamination assessment.

\paragraph{Future work.}
The Open LLM Leaderboard v1 provides per-item binary responses for 13K+ models---a unique resource for retrospective contamination audit once neural IRT estimation methods yield reliable discrimination parameters at scale.
Current classical IRT estimation (MML-EM) on leaderboard data produces discrimination parameters well below quality thresholds ($a$-parameter $r{=}0.13$, $N{=}178$ models), but advances in neural IRT estimation may close this gap.

%--------------------------------------------------------------
\section{Conclusion}
\label{sec:conclusion}

We introduce \textsc{IRT-FIT}, a contamination detection framework operating at the minimal-information level---requiring only binary response patterns.
Controlled experiments across three model families demonstrate detection with $|d|{>}2.0$ in two families (Qwen, Llama) at 15\% MMLU contamination, with bootstrap CIs quantifying estimation precision.
The $\theta$-shift mechanism provides mechanistic insight into detection boundary conditions.
While \textsc{IRT-FIT} incurs a 56\% effect-size cost relative to privileged methods, it requires zero model access and shows preliminary post-RL signal retention where token-probability methods fail.
By establishing the minimal-information operating point on the detection frontier, this work provides both a practical tool for closed-source model evaluation and a principled framework for understanding what binary response patterns can and cannot reveal about benchmark contamination.

%--------------------------------------------------------------
\bibliographystyle{plainnat}
\begin{thebibliography}{20}

\bibitem[Brown et~al.(2020)]{brown2020language}
Brown, T.~B., et~al.
\newblock Language models are few-shot learners.
\newblock In \emph{NeurIPS}, 2020.

\bibitem[Deng et~al.(2024)]{deng2024investigating}
Deng, C., et~al.
\newblock Investigating data contamination in modern benchmarks for large language models.
\newblock In \emph{ICLR}, 2024.

\bibitem[Dodge et~al.(2021)]{dodge2021documenting}
Dodge, J., et~al.
\newblock Documenting large webtext corpora: A case study on the {C}olossal {C}lean {C}rawled {C}orpus.
\newblock In \emph{EMNLP}, 2021.

\bibitem[Drasgow et~al.(1985)]{drasgow1985appropriateness}
Drasgow, F., Levine, M.~V., and Williams, E.~A.
\newblock Appropriateness measurement with polychotomous item response models and standardized indices.
\newblock \emph{British Journal of Mathematical and Statistical Psychology}, 38(1):67--86, 1985.

\bibitem[Hong et~al.(2021)]{hong2021multiscale}
Hong, M., Lin, L., and Cheng, Y.
\newblock Multiscale person-fit statistics.
\newblock \emph{Psychometrika}, 86:577--600, 2021.

\bibitem[Oren et~al.(2024)]{oren2024proving}
Oren, Y., et~al.
\newblock Proving test set contamination in black box language models.
\newblock In \emph{NeurIPS}, 2024.

\bibitem[PSN-IRT(2026)]{psnirt2026}
{PSN-IRT Authors}.
\newblock {PSN-IRT}: Psychometric scaling networks for item response theory.
\newblock In \emph{AAAI}, 2026.

\bibitem[Shi et~al.(2024)]{shi2024detecting}
Shi, W., et~al.
\newblock Detecting pretraining data from large language models.
\newblock In \emph{ICLR}, 2024.

\bibitem[Snijders(2001)]{snijders2001}
Snijders, T.~A.~B.
\newblock Asymptotic null distribution of person fit statistics with estimated person parameter.
\newblock \emph{Psychometrika}, 66(3):331--342, 2001.

\bibitem[Wang et~al.(2025)]{wang2025rlcontam}
Wang, X., et~al.
\newblock Benchmark contamination under post-reinforcement learning.
\newblock In \emph{ICLR}, 2026.

\bibitem[Zhang et~al.(2024)]{zhang2024min}
Zhang, Y., et~al.
\newblock {Min-K\%++}: Improved baseline for detecting pre-training data from large language models.
\newblock In \emph{ICLR}, 2024.

\bibitem[CoDeC(2026)]{codec2026}
{CoDeC Authors}.
\newblock {CoDeC}: Detecting contamination through in-context learning effects.
\newblock In \emph{ICLR}, 2026.

\bibitem[Meijer and Sijtsma(2001)]{meijer2001methodology}
Meijer, R.~R. and Sijtsma, K.
\newblock Methodology review: Evaluating person fit.
\newblock \emph{Applied Psychological Measurement}, 25(2):107--135, 2001.

\bibitem[IRT-LLM(2024)]{irt_llm_2024}
{IRT-LLM Authors}.
\newblock Item response theory for {LLM} evaluation.
\newblock 2024.

\end{thebibliography}

\end{document}
